\chapter*{Appendix}
\section*{Matlab code for the logistic map}
\label{sec:matlab_cobweb}
I took the liberty to use a slightly different code from the professor for that simulation:
\begin{lstlisting}[style=Matlab-editor, caption={Logistic Map Script}]
clear all

close all

set(groot, 'defaultFigureColor', 'w', ...
    'defaultAxesColor', 'w', ...
    'defaultAxesXColor', 'k', ...
    'defaultAxesYColor', 'k', ...
    'defaultTextColor', 'k');

N=200; % number of iterations

r=2.6;

x0=0.95;


xm=linspace(0,1,100);

ym=map(r,xm);

x = zeros(1,N);

xp=zeros(1,2*N+1);

yp=zeros(1,2*N+1);

x(1)=x0; % initial condition (can be anything from 0 to 1)

xp(1)=x0;


for n = 1:N

    x(n+1)=map(r,x(n));

    xp(2*n)=xp(2*n-1);

    yp(2*n)=x(n+1);

    xp(2*n+1)=x(n+1);

    yp(2*n+1)=x(n+1);

end

% Define the double-time map function
map2 = @(r, s) map(r, map(r, s));
ym2 = map2(r, xm);

% Compute trajectory and cobweb points for the double-time map
x2 = zeros(1, N);
xp2 = zeros(1, 2*N+1);
yp2 = zeros(1, 2*N+1);
x2(1) = x0;
xp2(1) = x0;

for n = 1:N
    x2(n+1) = map2(r, x2(n));
    xp2(2*n) = xp2(2*n-1);
    yp2(2*n) = x2(n+1);
    xp2(2*n+1) = x2(n+1);
    yp2(2*n+1) = x2(n+1);
end

% Create a single figure window with subplots
figure;

% First panel: Map function and Cobweb diagram
subplot(1, 2, 1);
plot(xm, xm, 'k--', 'LineWidth', 1);     % Identity line y = x
hold on;
plot(xm, ym, 'r', 'LineWidth', 1.5);        % Single-time map curve
plot(xp, yp, 'b');                           % Single map cobweb trajectory
title('T(x)');
xlabel('x_n');
ylabel('x_{n+1}');
legend('y = x', 'Single Map', 'Cobweb 1x');
grid on;

%Second panel: Double map
%subplot(1, 3, 2);
%plot(xm, xm, 'k--', 'LineWidth', 1);     % Identity line y = x
%hold on;
%plot(xm, ym2, 'm', 'LineWidth', 1.5);       % Double-time map curve
%plot(xp2, yp2, 'g');                       % Double map cobweb trajectory
%title('T^2(x)');
%xlabel('x_n');
%ylabel('x_{n+1}');
%legend('y = x', 'Double Map', 'Cobweb 2x');
%grid on;

% Second/third panel: Time series trajectory x(n)
subplot(1, 2, 2);
plot(x, '.-');
xlim([0 200]);
title('Trajectory x(n)');
xlabel('n');
ylabel('x_n');
grid on;

function m=map(r,s)
m=r*s.*(1-s);
%    m0=r*s.*(1-s);   % this is the double-time map
%    m=r*m0.*(1-m0);
end

\end{lstlisting}

One quick note: to display also the squared map, one must uncomment the middle panel and then change all the subplot() have 3 as the middle number, and also as the last one for the trajectory panel.